\section{Literature and original content}
We compared primary papers on complementarity and final-state boundaries, observer-relative states, horizon decoherence and gravitational localization, correlations between worldlines, matrix membrane indices and scattering, clock conditioning, noncommutative geometry, flux ensembles and environment-assisted correction. Searches combined the concepts of spectral flow, relative charge, horizon, conditional interior, irreversible records, interacting histories and matrix dynamics, and followed the closest papers' technical terms. For observational claims we used the original observational study or the publicly released table. The bibliography in the article and \texttt{evidence/primary\_sources.json} identify the specific input used.

The nearest conceptual precedents differ at identifiable steps. Horizon decoherence lets long-range fields carry a record; recovering an ordinary environment can change which phases are available. Measurement and feedback can source a classical gravitational field. Correlated-worldline gravity couples paths, while a final-state condition specifies a quantum boundary at a future singularity. Unitarily evolving membranes already exchange longitudinal momentum. CAT instead postulates a trace-preserving, charge-selecting contact in the causal past of the exterior boundary and tests the physical dimension of its inward section space. We found no realization of this complete prescription in the primary works compared. The observations and algebra identify its differences under specified preparations; a bibliographic search cannot guarantee worldwide priority.

\section{Coverage of the electronic evidence}
The table \texttt{evidence/result\_summaries.tsv} gives an individually written English account of every numbered result used for this release. The underlying selection criterion was introduced at the 21 November 2026, 18:54 JST research revision. The reference to a numbered result is an identifier for the included evidence. No preceding manuscript is needed.

\begin{center}\small
\begin{tabularx}{\linewidth}{@{}L{.20\linewidth}X@{}}\toprule
Result identifiers & How the calculations enter the argument\\\midrule
0093--0109 & Continuous Born division, material references, membrane forces, candidate mediators, causal onset, direct contacts, recoil and finite lifetimes. Together these motivate a contact with no ordinary interlayer tensor-wave channel.\\
0110--0123 & Correlation transfer, macroscopically different observer receipts, grains and human masses, time reversal and optional cosmological or galactic stress. Each has a material or spectral preparation condition.\\
0124--0141 & Finite deadlines, signed charge, local and returning radiation, winding, captured modes, phase-sensitive matching and terminal limitations. These give criteria for choosing a parent before Planck curvature.\\
0142--0151 & Construction of the parent, spin readout, all-species recovery, material clocks, ordinary vacuum waves, finite angular resolution, differential tests and existence of a classical landscape.\\
0152--0167 & Curvature, timelike and fermionic receivers, recoil, high-curvature continuation and stability, private overlaps, longitudinal flux and finite-time elastic response.\\
0168--0181 & Specified flux counts, observational branch tests, magnetized matter, radius dynamics, localized charge, long-range atomic forces and the comparison to ordinary IKKT and BFSS.\\
0182 & A signed clock reduction and one prepared source linking causal reception, recoil, strict sector closure and complete-history decoding.\\
0183 & Physical scalar sign, observer-specific inward fibers, shared records and a return experiment with two-sided readout calibration.\\
\bottomrule\end{tabularx}
\end{center}
The continuous charges, direct potentials and mediators explored in the early records are comparison candidates. The adopted spherical boundary uses integer graded bundles. The 60 percent figure in RES-0101 is a relative action-approximation error. The record includes the failed preparations. Examples include insufficient tidal excitation, a common-port support floor, radiation that has the right spectrum but loses a phase, and radii stabilized in a truncation that conflict with either perturbative stability or atomic forces. Numerical programs and formal modules are retained with their inputs. Detailed historical calculations and method notes are in Japanese; the English result guide and reviewer instructions identify what to run.

